\documentclass[10pt,a4paper]{article}
\usepackage{amsmath,amssymb,amsthm,graphicx}
\usepackage[no-math]{fontspec}
\newenvironment{authorclaim}[1]{\par\medskip\noindent\textbf{#1.}\ }{\par\medskip}
\allowdisplaybreaks
\title{Polynomial-time nilpotency-class reduction for exponent-p matrix-group isomorphism}
\author{Joshua A. Grochow, Youming Qiao}
\date{doi:10.4230/LIPIcs.CCC.2021.16; CC BY 4.0}
\begin{document}
\maketitle
\noindent\textit{Editorial scope.} Counted once: formal Theorem P nilpotency-class reduction. Preserves odd prime, exponent, class, field and matrix-generator input conditions. Full published local constructive correspondence argument and final author proof retained; imported tensor completeness and classical Lazard/BCH theory remain cited background. Three-generated extension, support proposition and runtime remarks are not extra tasks.\par\medskip
\noindent\textit{Allowed standard prerequisites (editorial).}\begin{itemize}
\item Lazard category correspondence, inverse BCH formulas, Sylow theory and imported tensor-isomorphism completeness reductions. Exact original Theorem 25, tensor/TI and Lie definitions, imported [28]/[32] reductions and full 74-entry bibliography retained. Source-local Proposition 26 matrix realization argument and entire Theorem P proof are included; imported theory not recursively reproved.
\end{itemize}\medskip
% Faithful official-publication transcription for TheoremP and its local proof, CCC2021.16.
\section*{Author statement: Theorem P}

\begin{authorclaim}{TheoremP}
Let $p$ be an odd prime. For groups given by generating sets of $m$ matrices of size $n\times n$ over $\mathbb F_{p^e}$, Group Isomorphism for $p$-groups of exponent $p$ and class $c<p$ reduces to Group Isomorphism for $p$-groups of exponent $p$ and class 2 in time $\operatorname{poly}(n,m,e\log p)$.
\end{authorclaim}

In fact, because the Lazard Correspondence works whenever all subgroups generated by 3 elements have nilpotency class $<p$, our reduction also works in this more general setting. For example, as a consequence of TheoremP, testing isomorphism of 5-groups in which every 3-generated subgroup has class 4 (the groups themselves may have larger class) reduces to testing isomorphism of 5-groups of class 2 in the matrix group model over fields of characteristic 5.
\begin{authorclaim}{Remark2}
Two additional results would suffice to get the analogous result in the Cayley table model. The first is to compute the Lazard Correspondence in the Cayley table model in time $\operatorname{poly}(|G|)$; we thank an anonymous ITCS reviewer for pointing out that this can be achieved by applying the matrix Lazard Correspondence (see Proposition26) to the left regular representation of the group on itself. The second is to improve the blow-up in the reduction from (Lie) Algebra Isomorphism to 3TI from [28]. Currently this reduction increases the dimension quadratically, which means the size of the group becomes $|G|^{O(\log|G|)}$ after the reduction; instead, we would need a reduction that increases the dimension only linearly.
\end{authorclaim}
\begin{authorclaim}{Remark3}
One may also ask whether our theorems can be combined, in order to get search-to-decision and counting-to-decision reductions for $p$-groups of class $c<p$ instead of only class 2. We believe this should be approachable, but again the quadratic increase in dimension in reductions, mentioned in the previous remark, gets in the way. The quadratic increase makes the square-root exponential reductions into ordinary exponential reductions, negating any gains.
\end{authorclaim}
\section*{Necessary author definitions (Sections2,3 and6.1)}
\paragraph{Vector spaces.} Let $\mathbb F$ be a field. In this paper we only consider finite-dimensional vector spaces over $\mathbb F$. We use $\mathbb F^n$ to denote the vector space of length-$n$ column vectors. The $i$th standard basis vector of $\mathbb F^n$ is denoted $\vec e_i$. Depending on the context, 0 may denote the zero vector space, a zero vector, or an all-zero matrix. For $S$ a set of vectors, we use $\langle S\rangle$ to denote the subspace spanned by elements in $S$.
\paragraph{Some groups.} The general linear group of degree $n$ over a field $\mathbb F$ is denoted by $\mathrm{GL}(n,\mathbb F)$.
\paragraph{Nilpotent groups.} If $A,B$ are two subsets of a group $G$, then $[A,B]$ denotes the subgroup generated by all elements of the form $[a,b]=aba^{-1}b^{-1}$, for $a\in A,b\in B$. The lower central series of a group $G$ is defined as follows: $\gamma_1(G)=G$, $\gamma_{k+1}(G)=[\gamma_k(G),G]$. A group is nilpotent if there is some $c$ such that $\gamma_{c+1}(G)=1$; the smallest such $c$ is called the nilpotency class of $G$, or sometimes just ``class'' when it is understood from context. A finite group is nilpotent if and only if it is the product of its Sylow subgroups; in particular, all groups of prime power order are nilpotent.
\paragraph{Matrices.} Let $\mathrm M(\ell\times n,\mathbb F)$ be the linear space of $\ell\times n$ matrices over $\mathbb F$, and $\mathrm M(n,\mathbb F):=\mathrm M(n\times n,\mathbb F)$. Given $A\in\mathrm M(\ell\times n,\mathbb F)$, $A^t$ denotes the transpose of $A$.
\begin{authorclaim}{Definition6 (The $d$-Tensor Isomorphism problem)}
$d$-Tensor Isomorphism over a field $\mathbb F$ is the problem: given two $d$-way arrays $\mathtt A=(a_{i_1,\ldots,i_d})$ and $\mathtt B=(b_{i_1,\ldots,i_d})$, where $i_k\in[n_k]$ for $k\in[d]$, and $a_{i_1,\ldots,i_d},b_{i_1,\ldots,i_d}\in\mathbb F$, decide whether there are $P_k\in\mathrm{GL}(n_k,\mathbb F)$ for $k\in[d]$, such that for all $i_1,\ldots,i_d$,
\[a_{i_1,\ldots,i_d}=\sum_{j_1,\ldots,j_d}b_{j_1,\ldots,j_d}(P_1)_{i_1,j_1}(P_2)_{i_2,j_2}\cdots(P_d)_{i_d,j_d}.\]
\end{authorclaim}
\paragraph{TI-completeness.} As the proof of TheoremP in Section6.3.1 uses a result on TI-completeness from [32], here we recall the definition of TI; see Definition6 for the $d$-Tensor Isomorphism problem.
\begin{authorclaim}{Definition22 ($d$TI, TI)}
For any field $\mathbb F$, $d\mathrm{TI}_{\mathbb F}$ denotes the class of problems that are polynomial-time Turing (Cook) reducible to $d$-Tensor Isomorphism over $\mathbb F$. Also let $\mathrm{TI}_{\mathbb F}=\bigcup_{d\ge1}d\mathrm{TI}_{\mathbb F}$.
\end{authorclaim}
The relationship between TI over different fields remains an intriguing open question [32], but here we will only need TI over $\mathbb F_p$. One of the the main results of [32] is that $\mathrm{TI}=d\mathrm{TI}$ for any fixed $d\ge3$.
\paragraph{Algebras and their algorithmic representations.} A Lie algebra $A$ consists of a vector space $V$ and a bilinear map $[,]:V\times V\to V$ that is alternating ($[v,v]=0$ for all $v\in V$; this is equivalent to skew-symmetry $[u,v]=-[v,u]$ in characteristic not 2) and satisfies the Jacobi identity $[x,[y,z]]+[z,[x,y]]+[y,[z,x]]=0$. The Jacobi identity is essentially the ``derivative'' of associativity.
After choosing an ordered basis $(b_1,\ldots,b_n)$ where $b_i\in\mathbb F^n$ of $V\cong\mathbb F^n$, this bilinear map $[,]$ can be represented by an $n\times n\times n$ 3-way array $\mathtt A$, such that $[b_i,b_j]=\sum_{k\in[n]}\mathtt A(i,j,k)b_k$. This is the structure constant representation of $A$. Algorithms for Lie algebras have been studied intensively in this model, e.g., [21,38].
It is also natural to consider matrix spaces that are closed under commutator. More specifically, let $A\le\mathrm M(n,\mathbb F)$ be a matrix space. If $A$ is closed under commutator, that is, for any $A,B\in A$, $[A,B]=AB-BA\in A$, then $A$ is a matrix Lie algebra with the product being the commutator. (Protip: one way to remember the Jacobi identity is to derive it as the natural identity among nested commutators of three matrices.) Algorithms for matrix Lie algebras have also been studied, e.g., [24,36,38].

\section*{6.3 Constructive Lazard's Correspondence and TheoremP}
The Lazard Correspondence [46] is a correspondence between certain classes of groups and Lie algebras, which extends the usual correspondence between Lie groups and Lie algebras (say, over $\mathbb R$) to some groups and Lie algebras in positive characteristic. Here we state just enough to give a sense of it; for further details and exposition we refer to Khukhro's book [43] and Naik's thesis [60]. While Naik's thesis is quite long, it also includes a reader's guide, and collects many results scattered across the literature or well-known to the experts in one place, building the theory from the ground up and with many examples.
Recall that a Lie ring is an abelian group $L$ equipped with a bilinear map $[,]$, called the Lie bracket, which is (1) alternating ($[x,x]=0$ for all $x\in L$) and (2) satisfies the Jacobi identity $[x,[y,z]]+[y,[z,x]]+[z,[x,y]]=0$ for all $x,y,z\in L$ (in some sense the ``derivative'' of the associativity equation). Let $L_1=L$, and $L_{i+1}=[L,L_i]$, which is the subgroup (of the underlying additive group) generated by all elements of the form $[x,y]$ for $x\in L,y\in L_i$. Then $L$ is nilpotent if $L_{c+1}=0$ for some finite $c$; the smallest such $c$ is the nilpotency class. (Lie algebras are just Lie rings over a field.)
The correspondence between Lie algebras and Lie groups over $\mathbb R$ uses the Baker--Campbell--Hausdorff (BCH) formula to convert between a Lie algebra and a Lie group, so we start there. The BCH formula is the solution to the problem that for non-commuting matrices $X,Y$, $e^Xe^Y\ne e^{X+Y}$ in general (where the matrix exponential here is defined using the power series for $e^x$). Rather, using commutators $[A,B]=AB-BA$, we have
\[\exp(X)\exp(Y)=\exp\left(X+Y+\frac12[X,Y]+\frac1{12}([X,[X,Y]]-[Y,[X,Y]])+\cdots\right),\]
where the remaining terms are iterated commutators that all involve at least 4 $X$s and $Y$s, and successive terms involve more and more. Applying the exponential function to a Lie algebra in characteristic zero yields a Lie group. The BCH formula can be inverted, giving the correspondence in the other direction.
In a nilpotent Lie algebra, the BCH formula has only finitely many nonzero terms, so issues of convergence disappear and we may consider applying the correspondence over finite fields or rings; the only remaining obstacle is that the denominators appearing in the formula must be units in the ring. It turns out that the correspondence continues to work in characteristic $p$ so long as one does not need to use the $p$-th term of the BCH formula (which includes division by $p$), and the latter is avoided whenever a nilpotent group has class strictly less than $p$, or even when all subgroups generated by at most 3 elements have class strictly less than $p$. While the correspondence does apply more generally, here we only state the version for finite groups. For any fixed nilpotency class $c$, computing the Lazard Correspondence is efficient in theory; for how to compute it in practice when the groups are given by polycyclic presentations, see [20].
Let $\mathrm{Grp}_{p,n,c}$ denote the set of finite groups of order $p^n$ and class $c$, and let $\mathrm{Lie}_{p,n,c}$ denote the set of Lie rings of order $p^n$ and class $c$. We note that for nilpotency class 2, the Baer Correspondence is the same as the Lazard Correspondence.
\begin{authorclaim}{Theorem25 (Lazard Correspondence for finite groups [46], see, e.g., [43, Ch.9 \&10] or [60, Ch.6])}
For any prime $p$ and any $1\le c<p$, there are functions $\log:\mathrm{Grp}_{p,n,c}\leftrightarrow\mathrm{Lie}_{p,n,c}:\exp$ such that (1) $\log$ and $\exp$ are inverses of one another, (2) two groups $G,H\in\mathrm{Grp}_{p,n,c}$ are isomorphic if and only if $\log(G)$ and $\log(H)$ are isomorphic, and (3) if $G$ has exponent $p$, then the underlying abelian group of $\log(G)$ has exponent $p$. More strongly, $\log$ is an isomorphism of categories $\mathrm{Grp}_{p,n,c}\cong\mathrm{Lie}_{p,n,c}$.
\end{authorclaim}
Part(3) can be found as a special case of [60, Lemma6.1.2].
For $p$-groups given by $d\times d$ matrices over the finite field $\mathbb F_{p^e}$, we will need one additional fact about the correspondence, namely that it also results in a Lie algebra of $d\times d$ matrices. (Being able to bound the dimension of the Lie algebra and work with it in a simple linear-algebraic way seems crucial for our reduction to work efficiently.) In fact, the BCH Correspondence is easier to see for matrix groups using the matrix exponential and matrix logarithm; most of the work for BCH and Lazard is to get the correspondence to work even without the matrices. In some sense, this is thus the ``original'' setting of this correspondence. Though it is surely not new, we could not find a convenient reference for this fact about matrix groups over finite fields, so we state it formally here.
\begin{authorclaim}{Proposition26 (cf. [43, Exercise10.6])}
Let $G\le\mathrm{GL}(d,\mathbb F_{p^e})$ be a finite $p$-subgroup of exponent $p$, consisting of $d\times d$ matrices over a finite field of characteristic $p$. Then $\log(G)$ (from the Lazard Correspondence) can be realized as a finite Lie subalgebra of $de\times de$ matrices over $\mathbb F_p$. Given a generating set for $G$ of $m$ matrices, a generating set for $\log(G)$ can be constructed in $\operatorname{poly}(d,n,e\log p)$ time.
\end{authorclaim}
Khukhro [43] gives the characteristic zero analogue of this result (minus the straightforward complexity analysis) for the full group of upper unitriangular matrices as Exercise10.6. One way to see Proposition26 is to use the characteristic zero result, apply the fact that these isomorphisms are in fact equivalence of categories (and thus hold for subgroups/subalgebras as well), and note that the same formulae in characteristic zero apply in characteristic $p$ so long as one never needs to divide by $p$. We now sketch the argument.
\paragraph{Proof sketch.}
First we use the standard embedding of $\mathrm{GL}(d,\mathbb F_{p^e})$ into $\mathrm{GL}(de,\mathbb F_p)$ (replace each element by an $e\times e$ block which is the left regular representation of $\mathbb F_{p^e}$ acting on itself as an $e$-dimensional $\mathbb F_p$-vector space), to realize $G$ as a subgroup of $\mathrm{GL}(de,\mathbb F_p)$. $G$ is conjugate in $\mathrm{GL}(de,\mathbb F_p)$ to a group of upper unitriangular matrices (upper triangular with all 1s on the diagonal); this is a standard fact that can be seen in several ways, for example, by noting that the group $U$ of all upper unitriangular matrices in $\mathrm{GL}(de,\mathbb F_p)$ is a Sylow $p$-subgroup, and applying Sylow's Theorem. (Note that we do not need to do this conjugation algorithmically, though it is possible to do so [27,36,64]; this is only for the proof.) Thus we may write every $g\in G$ as $1+n$, where the sum here is the ordinary sum of matrices, 1 denotes the identity matrix, and $n$ is strictly upper triangular. To see that we can truncate the Taylor series for logarithm before the $p$-th term (thus avoiding needing to divide by $p$), note that $(1+n)^p=1$ since $G$ is exponent $p$. We have $(1+n)^p=1^p+\binom p1n+\binom p2n^2+\cdots+\binom p{p-1}n^{p-1}+n^p$. Since these are matrices over a field of characteristic $p$, and $p\mid\binom pi$ for all $1\le i\le p-1$, all the intermediate terms vanish and we have that $(1+n)^p=1^p+n^p$. Thus $1=(1+n)^p=1+n^p$, so we get that $n^p=0$. Thus, in the the Taylor series for the logarithm
\[\log(1+n)=n-\frac{n^2}2+\frac{n^3}3-\cdots\]
the last nonzero term is $n^{p-1}/(p-1)$, so we may use this Taylor series even over $\mathbb F_{p^e}$.
The main things to check are that the set $\log(G):=\{\log(1+n):1+n\in G\}$ is closed under scalar multiplication, matrix addition, and matrix commutator $[X,Y]=XY-YX$. Suppose $g_1,g_2$ are matrices in $G$, and write them as $g_i=1+n_i$ ($i=1,2$), as above. We recall that, because $n_i^p=0$ from above, the power series for both $\log$ and $\exp$ work to compute the matrix logarithm and exponential over $\mathbb F_{p^e}$, respectively, and that the usual rules of logarithms are satisfied for a single matrix $A$: whenever $A\in\mathrm M_{de}(\mathbb F_p)$ satisfies $A^p=0$, we have $\log\exp A=A$, $\exp\log(1+A)=1+A$, $\exp(nA)=(\exp A)^n$ for $n\in\mathbb Z$, and $\log((1+A)^n)=n\log(1+A)$.
\begin{itemize}
\item Scalar multiplication: For $\alpha\in\mathbb F_p$, we show that $\alpha\log(1+n_1)$ is in $\log(G)$. This is easy to show, as it follows directly from the rules of logarithms just mentioned: $\alpha\log(1+n_1)=\log((1+n_1)^\alpha)$ where on the right-hand side we treat $\alpha$ as an integer in the range $[0,p-1]$. (This is the only point where we are using that we are working over $\mathbb F_p$ now rather than $\mathbb F_{p^e}$.)
\item Addition: Let $x_i=\log(1+n_i)$ for $i=1,2$. We want to show that $x_1+x_2$ is in $\log(G)$, or equivalently that $\exp(x_1+x_2)\in G$. This follows from the first inverse BCH formula $h_1$, which satisfies $\exp(\hat x_1+\hat x_2)=h_1(\exp(\hat x_1),\exp(\hat x_2))$ for $\hat x_i$ in the free nilpotent-of-class-$c$ $\mathbb F_p$-Lie algebra, and then we may apply the homomorphism from the latter algebra to the subalgebra of $\mathrm M_n(\mathbb F_p)$ generated by the $n_i$ to see that the same formula works. (We note, because a reviewer asked, that here we do not need this entire subalgebra to be in $\{g-1:g\in G\}$; the use of that subalgebra is just convenient for talking about algebra homomorphisms in the proof. Rather, it suffices that the preceding equation holds for these particular elements $n_i$, which are by definition of the form $g_i-1$ for some matrices $g_i\in G$.)
\item Commutator: $[\log(1+n_1),\log(1+n_2)]$. A similar argument as in the previous case works, using the second inverse BCH formula $h_2$, which satisfies $\exp([\hat x_1,\hat x_2])=h_2(\exp(\hat x_1),\exp(\hat x_2))$.
\end{itemize}
Equivalently, we may note that the derivation of the inverse BCH formulas in [43] uses a free nilpotent associative algebra as an ambient setting in which both the group (or rather, $n$ such that $1+n$ is in the group) and the corresponding Lie algebra live; in our case, we may replace the ambient free nilpotent associative algebra with the algebra of $de\times de$ strictly upper-triangular matrices over $\mathbb F_p$, and all the derivations remain the same, mutatis mutandis. See, for example, [43,p.105, ``Another remark...'']. \hfill$\square$
\subsection*{6.3.1 Class reduction in $p$-group isomorphism testing}
Proposition26 now allows us to prove TheoremP.
\begin{proof}[Proof of TheoremP]
By the Lazard Correspondence (reproduced as Theorem25) two $p$-groups of exponent $p$ and class $c<p$ are isomorphic if and only if their corresponding $\mathbb F_p$-Lie algebras are. By Proposition26, we can construct a generating set for the corresponding $\mathbb F_p$-Lie algebra by applying the power series for logarithm to the generating matrices of $G$. This Lie algebra is thus a subalgebra of $ne\times ne$ matrices over $\mathbb F_p$, so we can generate a basis for the entire Lie algebra (using the linear-algebra version of breadth-first search; its dimension is $\le(ne)^2$) and compute its structure constants in time polynomial in $n,m$, and $e\log p$. Then use [28] to reduce isomorphism of Lie algebras to 3-Tensor Isomorphism, and then use the fact that isomorphism of $p$-groups of exponent $p$ and class 2 given by a matrix generating set over $\mathbb F_p$ is TI-complete [32] to reduce to the latter problem.
\end{proof}

\begin{thebibliography}{74}

\bibitem[1]{ref1}
Eric Allender and Bireswar Das. Zero knowledge and circuit minimization. Inf. Comput., 256:2–8, 2017. doi:10.1016/j.ic.2017.04.004.

\bibitem[2]{ref2}
Vikraman Arvind and Jacobo Torán. Isomorphism testing: Perspective and open problems. Bulletin of the EATCS, 86:66–84, 2005.

\bibitem[3]{ref3}
László Babai. Graph isomorphism in quasipolynomial time [extended abstract]. In Proceedings of the 48th Annual ACM SIGACT Symposium on Theory of Computing, STOC 2016, pages 684–697, 2016. arXiv:1512.03547 [cs.DS] version 2. doi:10.1145/2897518.2897542.

\bibitem[4]{ref4}
László Babai, Paolo Codenotti, Joshua A. Grochow, and Youming Qiao. Code equivalence and group isomorphism. In Proceedings of the Twenty-Second Annual ACM–SIAM Symposium on Discrete Algorithms (SODA11), pages 1395–1408, Philadelphia, PA, 2011. SIAM. doi: 10.1137/1.9781611973082.107.

\bibitem[5]{ref5}
László Babai, Paolo Codenotti, and Youming Qiao. Polynomial-time isomorphism test for groups with no abelian normal subgroups - (extended abstract). In Automata, Languages, and Programming - 39th International Colloquium, ICALP 2012, Proceedings, Part I, pages 51–62, 2012. doi:10.1007/978-3-642-31594-7\_5.

\bibitem[6]{ref6}
László Babai and Youming Qiao. Polynomial-time isomorphism test for groups with Abelian Sylow towers. In 29th STACS, pages 453–464. Springer LNCS 6651, 2012. doi:10.4230/ LIPIcs.STACS.2012.453.

\bibitem[7]{ref7}
Reinhold Baer. Groups with abelian central quotient group. Trans. AMS, 44(3):357–386, 1938. doi:10.1090/S0002-9947-1938-1501972-1.

\bibitem[8]{ref8}
Mihir Bellare and Shafi Goldwasser. The complexity of decision versus search. SIAM J. Comput., 23(1):97–119, 1994. doi:10.1137/S0097539792228289.

\bibitem[9]{ref9}
Hans Ulrich Besche and Bettina Eick. Construction of finite groups. J. Symb. Comput., 27(4):387–404, 1999. doi:10.1006/jsco.1998.0258.

\bibitem[10]{ref10}
Hans Ulrich Besche, Bettina Eick, and E.A. O’Brien. A millennium project: Constructing small groups. Intern. J. Alg. and Comput, 12:623–644, 2002. doi:10.1142/S0218196702001115.

\bibitem[11]{ref11}
Anton Betten, Michael Braun, Harald Fripertinger, Adalbert Kerber, Axel Kohnert, and Alfred Wassermann. Error-correcting linear codes: Classification by isometry and applications, volume 18. Springer Science and Business Media, 2006.

\bibitem[12]{ref12}
Peter A. Brooksbank, Joshua A. Grochow, Yinan Li, Youming Qiao, and James B. Wilson. Incorporating Weisfeiler–Leman into algorithms for group isomorphism. arXiv:1905.02518 [cs.CC], 2019.

\bibitem[13]{ref13}
Peter A. Brooksbank, Yinan Li, Youming Qiao, and James B. Wilson. Improved algorithms for alternating matrix space isometry: From theory to practice. In Fabrizio Grandoni, Grzegorz Herman, and Peter Sanders, editors, 28th Annual European Symposium on Algorithms, ESA 2020, September 7-9, 2020, Pisa, Italy (Virtual Conference), volume 173 of LIPIcs, pages 26:1–26:15. Schloss Dagstuhl - Leibniz-Zentrum für Informatik, 2020. doi:10.4230/LIPIcs. ESA.2020.26.

\bibitem[14]{ref14}
Peter A. Brooksbank and Eugene M. Luks. Testing isomorphism of modules. J. Algebra, 320(11):4020–4029, 2008. doi:10.1016/j.jalgebra.2008.07.014.

\bibitem[15]{ref15}
Peter A. Brooksbank, Joshua Maglione, and James B. Wilson. A fast isomorphism test for groups whose Lie algebra has genus 2. J. Algebra, 473:545–590, 2017. doi:10.1016/j. jalgebra.2016.12.007.

\bibitem[16]{ref16}
Peter A. Brooksbank and E. A. O’Brien. Constructing the group preserving a system of forms. Internat. J. Algebra Comput., 18(2):227–241, 2008. doi:10.1142/S021819670800441X.

\bibitem[17]{ref17}
John J. Cannon and Derek F. Holt. Automorphism group computation and isomorphism testing in finite groups. J. Symbolic Comput., 35(3):241–267, 2003. doi:10.1016/S0747-7171(02) 00133-5.

\bibitem[18]{ref18}
Xi Chen, Xiaotie Deng, and Shang-Hua Teng. Settling the complexity of computing two-player Nash equilibria. J. ACM, 56(3):Art. 14, 57, 2009. doi:10.1145/1516512.1516516.

\bibitem[19]{ref19}
Alexander Chistov, Gábor Ivanyos, and Marek Karpinski. Polynomial time algorithms for modules over finite dimensional algebras. In Proceedings of the 1997 International Symposium on Symbolic and Algebraic Computation, ISSAC ’97, pages 68–74. ACM, 1997. doi:10.1145/ 258726.258751.

\bibitem[20]{ref20}
Serena Cicalò, Willem A. de Graaf, and Michael Vaughan-Lee. An effective version of the Lazard correspondence. J. Algebra, 352(1):430–450, 2012. doi:10.1016/j.jalgebra.2011.11.031.

\bibitem[21]{ref21}
W.A. de Graaf. Lie Algebras: Theory and Algorithms, volume 56 of North-Holland Mathemat- ical Library. Elsevier Science, 2000.

\bibitem[22]{ref22}
Holger Dell, Thore Husfeldt, Dániel Marx, Nina Taslaman, and Martin Wahlén. Exponential time complexity of the permanent and the Tutte polynomial. ACM Trans. Algorithms, 10(4):Art. 21, 32, 2014. doi:10.1145/2635812.

\bibitem[23]{ref23}
Sean Eberhard and Stefan-C. Virchow. Random generation of the special linear group. Transactions of the American Mathematical Society, page 1, 2020. doi:10.1090/tran/8009.

\bibitem[24]{ref24}
Wayne Eberly and Mark Giesbrecht. Efficient decomposition of associative algebras over finite fields. Journal of Symbolic Computation, 29(3):441–458, 2000. doi:10.1006/jsco.1999.0308.

\bibitem[25]{ref25}
Bettina Eick, C. R. Leedham-Green, and E. A. O’Brien. Constructing automorphism groups of p-groups. Comm. Algebra, 30(5):2271–2295, 2002. doi:10.1081/AGB-120003468.

\bibitem[26]{ref26}
V. Felsch and J. Neubüser. On a programme for the determination of the automorphism group of a finite group. In Pergamon J. Leech, editor, Computational Problems in Abstract Algebra (Proceedings of a Conference on Computational Problems in Algebra, Oxford, 1967), pages 59–60, Oxford, 1970.

\bibitem[27]{ref27}
Katalin Friedl and Lajos Rónyai. Polynomial time solutions of some problems in computational algebra. In Robert Sedgewick, editor, Proceedings of the 17th Annual ACM Symposium on Theory of Computing, May 6-8, 1985, Providence, Rhode Island, USA, pages 153–162. ACM, 1985. doi:10.1145/22145.22162.

\bibitem[28]{ref28}
Vyacheslav Futorny, Joshua A. Grochow, and Vladimir V. Sergeichuk. Wildness for tensors. Lin. Alg. Appl., 566:212–244, 2019. doi:10.1016/j.laa.2018.12.022.

\bibitem[29]{ref29}
Joshua A. Grochow. Answer to “what is the hardest instance for the group isomorph- ism problem?”. Theoretical Computer Science Stack Exchange. URL: https://cstheory. stackexchange.com/a/42551/129.

\bibitem[30]{ref30}
Joshua A. Grochow and Youming Qiao. Polynomial-time isomorphism test of groups that are tame extensions - (extended abstract). In Algorithms and Computation - 26th International Symposium, ISAAC 2015, Nagoya, Japan, December 9-11, 2015, Proceedings, pages 578–589, 2015. doi:10.1007/978-3-662-48971-0\_49.

\bibitem[31]{ref31}
Joshua A. Grochow and Youming Qiao. Algorithms for group isomorphism via group extensions and cohomology. SIAM J. Comput., 46(4):1153–1216, 2017. Preliminary version in IEEE Conference on Computational Complexity (CCC) 2014 (DOI:10.1109/CCC.2014.19). Also available as arXiv:1309.1776 [cs.DS] and ECCC Technical Report TR13-123. doi:10.1137/ 15M1009767.

\bibitem[32]{ref32}
Joshua A. Grochow and Youming Qiao. On the complexity of isomorphism problems for tensors, groups, and polynomials I: Tensor Isomorphism-completeness. In ITCS, page to appear, 2021. arXiv:1907.00309.

\bibitem[33]{ref33}
Martin Grohe and Pascal Schweitzer. The graph isomorphism problem. Commun. ACM, 63(11):128–134, 2020. doi:10.1145/3372123.

\bibitem[34]{ref34}
Xiaoyu He and Youming Qiao. On the Baer–Lovász–Tutte construction of groups from graphs: isomorphism types and homomorphism notions. arXiv:2003.07200 [math.CO], 2020.

\bibitem[35]{ref35}
Russell Impagliazzo and Ramamohan Paturi. On the complexity of k-SAT. J. Comput. System Sci., 62(2):367–375, 2001. Special issue on the Fourteenth Annual IEEE Conference on Computational Complexity (Atlanta, GA, 1999). doi:10.1006/jcss.2000.1727.

\bibitem[36]{ref36}
Gábor Ivanyos. Fast randomized algorithms for the structure of matrix algebras over finite fields. In Proceedings of the 2000 international symposium on Symbolic and algebraic computation, pages 175–183. ACM, 2000. doi:10.1145/345542.345620.

\bibitem[37]{ref37}
Gábor Ivanyos, Marek Karpinski, and Nitin Saxena. Deterministic polynomial time algorithms for matrix completion problems. SIAM J. Comput., 39(8):3736–3751, 2010. doi:10.1137/ 090781231.

\bibitem[38]{ref38}
Gábor Ivanyos and Lajos Rónyai. Computations in associative and Lie algebras. In Some tapas of computer algebra, pages 91–120. Springer, 1999. doi:10.1007/978-3-662-03891-8\_5.

\bibitem[39]{ref39}
Zhengfeng Ji, Youming Qiao, Fang Song, and Aaram Yun. General linear group action on tensors: A candidate for post-quantum cryptography. In Dennis Hofheinz and Alon Rosen, editors, Theory of Cryptography - 17th International Conference, TCC 2019, Nuremberg, Germany, December 1-5, 2019, Proceedings, Part I, volume 11891 of Lecture Notes in Computer Science, pages 251–281. Springer, 2019. Preprint arXiv:1906.04330 [cs.CR]. doi:10.1007/ 978-3-030-36030-6\_11.

\bibitem[40]{ref40}
William M. Kantor. Some topics in asymptotic group theory. Groups, Combinatorics and Geometry (Durham, pages 403–421, 1990.

\bibitem[41]{ref41}
William M Kantor and Alexander Lubotzky. The probability of generating a finite classical group. Geometriae Dedicata, 36(1):67–87, 1990.

\bibitem[42]{ref42}
Neeraj Kayal and Timur Nezhmetdinov. Factoring groups efficiently. In Susanne Albers, Alberto Marchetti-Spaccamela, Yossi Matias, Sotiris E. Nikoletseas, and Wolfgang Thomas, editors, Automata, Languages and Programming, 36th International Colloquium, ICALP 2009, Rhodes, Greece, July 5-12, 2009, Proceedings, Part I, volume 5555 of Lecture Notes in Computer Science, pages 585–596. Springer, 2009. Preprint ECCC Tech. Report TR08-074. doi:10.1007/978-3-642-02927-1\_49.

\bibitem[43]{ref43}
E. I. Khukhro. p-automorphisms of finite p-groups, volume 246 of London Mathematical Society Lecture Note Series. Cambridge University Press, Cambridge, 1998. doi:10.1017/ CBO9780511526008.

\bibitem[44]{ref44}
Johannes Köbler, Uwe Schöning, and Jacobo Torán. The graph isomorphism problem: its structural complexity. Birkhauser Verlag, Basel, Switzerland, Switzerland, 1993. doi:10.1007/ 978-1-4612-0333-9.

\bibitem[45]{ref45}
Tamara G Kolda and Brett W Bader. Tensor decompositions and applications. SIAM review, 51(3):455–500, 2009. doi:10.1137/07070111X.

\bibitem[46]{ref46}
Michel Lazard. Sur les groupes nilpotents et les anneaux de Lie. Ann. Sci. Ecole Norm. Sup. (3), 71:101–190, 1954. doi:10.24033/asens.1021.

\bibitem[47]{ref47}
François Le Gall. Efficient isomorphism testing for a class of group extensions. In Proc. 26th STACS, pages 625–636, 2009. doi:10.4230/LIPIcs.STACS.2009.1830.

\bibitem[48]{ref48}
Mark L. Lewis and James B. Wilson. Isomorphism in expanding families of indistinguishable groups. Groups Complex. Cryptol., 4(1):73–110, 2012. doi:10.1515/gcc-2012-0008.

\bibitem[49]{ref49}
Yinan Li and Youming Qiao. Linear algebraic analogues of the graph isomorphism problem and the Erdős–Rényi model. In Chris Umans, editor, 58th IEEE Annual Symposium on Foundations of Computer Science, FOCS 2017, pages 463–474. IEEE Computer Society, 2017. doi:10.1109/FOCS.2017.49.

\bibitem[50]{ref50}
Richard J. Lipton, Lawrence Snyder, and Yechezkel Zalcstein. The complexity of word and isomorphism problems for finite groups. Yale University Department of Computer Science Research Report \# 91, 1977. URL: https://apps.dtic.mil/dtic/tr/fulltext/u2/a053246. pdf.

\bibitem[51]{ref51}
Eugene M. Luks. Computing in solvable matrix groups. In FOCS 1992, 33rd Annual Symposium on Foundations of Computer Science, pages 111–120. IEEE Computer Society, 1992. doi:10.1109/SFCS.1992.267813.

\bibitem[52]{ref52}
Eugene M. Luks. Permutation groups and polynomial-time computation. In Groups and computation (New Brunswick, NJ, 1991), volume 11 of DIMACS Ser. Discrete Math. Theoret. Comput. Sci., pages 139–175. Amer. Math. Soc., Providence, RI, 1993.

\bibitem[53]{ref53}
Eugene M. Luks. Hypergraph isomorphism and structural equivalence of boolean functions. In Proceedings of the Thirty-First Annual ACM Symposium on Theory of Computing, May 1-4, 1999, Atlanta, Georgia, USA, pages 652–658, 1999. doi:10.1145/301250.301427.

\bibitem[54]{ref54}
Rudolf Mathon. A note on the graph isomorphism counting problem. Information Processing Letters, 8(3):131–136, 1979.

\bibitem[55]{ref55}
Brendan D. McKay. Practical graph isomorphism. Congr. Numer., pages 45–87, 1980.

\bibitem[56]{ref56}
Brendan D. McKay and Adolfo Piperno. Practical graph isomorphism, II. Journal of Symbolic Computation, 60(0):94–112, 2014. doi:10.1016/j.jsc.2013.09.003.

\bibitem[57]{ref57}
Alan H. Mekler. Stability of nilpotent groups of class 2 and prime exponent. The Journal of Symbolic Logic, 46(4):781–788, 1981.

\bibitem[58]{ref58}
Gary L. Miller. On the $n^{\log n}$ isomorphism technique (a preliminary report). In STOC, pages 51–58. ACM, 1978. doi:10.1145/800133.804331.

\bibitem[59]{ref59}
Takunari Miyazaki. Luks’s reduction of graph isomorphism to code equivalence. Comment to E. W. Clark, https://groups.google.com/forum/\#!msg/sci.math.research/puZxGj9HXKI/ CeyH2yyyNFUJ, 1996.

\bibitem[60]{ref60}
Vipul Naik. Lazard correspondence up to isoclinism. PhD thesis, The University of Chicago, 2013. URL: https://vipulnaik.com/thesis/.

\bibitem[61]{ref61}
Jacques Patarin. Hidden fields equations (HFE) and isomorphisms of polynomials (IP): two new families of asymmetric algorithms. In Advances in Cryptology - EUROCRYPT ’96, Inter- national Conference on the Theory and Application of Cryptographic Techniques, Saragossa, Spain, May 12-16, 1996, Proceeding, pages 33–48, 1996. doi:10.1007/3-540-68339-9\_4.

\bibitem[62]{ref62}
Erez Petrank and Ron M. Roth. Is code equivalence easy to decide? IEEE Trans. Inf. Theory, 43(5):1602–1604, 1997. doi:10.1109/18.623157.

\bibitem[63]{ref63}
Youming Qiao, Jayalal M. N. Sarma, and Bangsheng Tang. On isomorphism testing of groups with normal Hall subgroups. In Proc. 28th STACS, pages 567–578, 2011. doi: 10.4230/LIPIcs.STACS.2011.567.

\bibitem[64]{ref64}
Lajos Rónyai. Computing the structure of finite algebras. J. Symb. Comput., 9(3):355–373, 1990. doi:10.1016/S0747-7171(08)80017-X.

\bibitem[65]{ref65}
David J. Rosenbaum. Bidirectional collision detection and faster deterministic isomorphism testing. arXiv preprint arXiv:1304.3935 [cs.DS], 2013.

\bibitem[66]{ref66}
David J. Rosenbaum. Breaking the $n^{\log n}$ barrier for solvable-group isomorphism. In Proceedings of the Twenty-Fourth Annual ACM-SIAM Symposium on Discrete Algorithms, pages 1054–1073. SIAM, 2013. Preprint arXiv:1205.0642 [cs.DS]. doi:10.1137/1.9781611973105.76.

\bibitem[67]{ref67}
Nicolas Sendrier and Dimitris E. Simos. The hardness of code equivalence over $\mathbb F_q$ and its application to code-based cryptography. In International Workshop on Post-Quantum Cryptography, pages 203–216. Springer, 2013.

\bibitem[68]{ref68}
Seinosuke Toda. PP is as hard as the polynomial-time hierarchy. SIAM J. Comput., 20(5):865– 877, 1991. doi:10.1137/0220053.

\bibitem[69]{ref69}
Leslie G. Valiant. Relative complexity of checking and evaluating. Information Processing Lett., 5(1):20–23, 1976/77. doi:10.1016/0020-0190(76)90097-1.

\bibitem[70]{ref70}
James Wilson. 2014 conference on Groups, Computation, and Geometry at Colorado State University, co-organized by P. Brooksbank, A. Hulpke, T. Penttila, J. Wilson, and W. Kantor. Personal communication, 2014.

\bibitem[71]{ref71}
James B. Wilson. Decomposing p-groups via Jordan algebras. J. Algebra, 322(8):2642–2679, 2009. doi:10.1016/j.jalgebra.2009.07.029.

\bibitem[72]{ref72}
James B. Wilson. Finding direct product decompositions in polynomial time. arXiv:1005.0548 [math.GR], 2010.

\bibitem[73]{ref73}
James B. Wilson. Existence, algorithms, and asymptotics of direct product decompositions, I. Groups Complex. Cryptol., 4(1):33–72, 2012. doi:10.1515/gcc-2012-0007.

\bibitem[74]{ref74}
V. N. Zemlyachenko, N. M. Korneenko, and R. I. Tyshkevich. Graph isomorphism problem. J. Soviet Math., 29(4):1426–1481, May 1985. doi:10.1007/BF02104746.
\end{thebibliography}
\end{document}
